% ============================================================
% SECCION 1 - INTRODUCCION (Robinson Jeanpierre / jean200525)
% Problema, contexto, objetivos y estructura del documento.
% ============================================================
\section{Introducci\'on}
\label{sec:introduccion}

\subsection{Problema y contexto}

La finca AgroMoreira es una unidad productiva familiar de 16 hect\'areas de policultivo ---palma, pl\'atano verde y cacao--- ubicada en la zona rural de Quevedo (Ecuador), con un equipo de siete trabajadores encabezado por un administrador con formaci\'on de cuarto nivel en biotecnolog\'ia agropecuaria. Pese a esa formaci\'on, la gesti\'on de la finca se sostiene sobre dos soportes fr\'agiles: registros manuales en bit\'acoras f\'isicas y comunicaci\'on verbal entre el administrador y los jornaleros.

Las consecuencias est\'an documentadas en la evidencia de campo del proyecto (EV-01 a EV-07): p\'erdida de control de producci\'on cuando no se registraron las fundas semanales de pl\'atano, proyecciones err\'oneas derivadas de datos incompletos, imposibilidad de reconstruir el historial de aplicaciones fitosanitarias y reacci\'on tard\'ia ante enfermedades como la moniliasis en cacao y la sigatoka en pl\'atano, a las que se suma una enfermedad nueva sin identificar que afecta la producci\'on. A ello se agrega una brecha estructural entre los perfiles de usuario: un administrador que decide con datos, prefiere gestionar desde el celular y pide expl\'icitamente alertas autom\'aticas; y jornaleros que ejecutan labores con instrucciones semanales verbales y muestran una disposici\'on neutra frente a la tecnolog\'ia.

El problema que este proyecto integrador aborda es, entonces, doble: (i)~la ausencia de un sistema centralizado de gesti\'on agr\'icola que capture la informaci\'on operativa de parcelas, cultivos, labores, cosechas e inventario; y (ii)~la ausencia de capacidad analítica para anticipar problemas fitosanitarios y de abastecimiento antes de que se traduzcan en p\'erdidas. La ingenier\'ia de requisitos es la disciplina que permite especificar ambos frentes con rigor verificable \cite{swebok2024}.

\subsection{Soluci\'on propuesta: el SGA}

El Sistema de Gesti\'on Agr\'icola con Inteligencia Artificial (SGA) se concibe como una aplicaci\'on web que centraliza la informaci\'on de la finca ---parcelas, cultivos de verde y cacao, trabajadores, labores, cosechas, insumos y herramientas--- y que incorpora dos componentes de inteligencia artificial: IA-01, recomendador de manejo agr\'icola basado en los registros del propio sistema, e IA-02, detector de plagas y enfermedades a partir de im\'agenes capturadas en campo. Ambos componentes se especifican con umbrales cuantitativos de rendimiento, equidad y explicabilidad, clasificaci\'on de riesgo conforme al Reglamento (UE) 2024/1689 y cumplimiento de la normativa ecuatoriana de protecci\'on de datos personales \cite{eu20241689,lopdp2021}.

Quedan deliberadamente fuera del alcance de esta versi\'on la aplicaci\'on m\'ovil nativa y los sensores IoT de campo, documentados como \emph{Won't Have} en la priorizaci\'on MoSCoW: el acceso web responsivo cubre las necesidades actuales declaradas por los stakeholders, incluida la preferencia del administrador por el navegador del celular.

\subsection{Objetivos del proyecto integrador}

\subsubsection*{Objetivo general}

Integrar, auditar y defender la Especificaci\'on de Requisitos de Software (ERS) del SGA de AgroMoreira como cierre del proyecto integrador de la asignatura, aplicando m\'etricas de calidad derivadas de ISO/IEC 25010:2023, verificando la trazabilidad end-to-end del requisito a la prueba y especificando los requisitos de los componentes de inteligencia artificial \cite{iso25010}.

\subsubsection*{Objetivos espec\'ificos}

\begin{enumerate}[leftmargin=0.9cm]
  \item Auditar la calidad del ERS final mediante seis m\'etricas operativas (M1--M6) calculadas sobre conteos reales, publicados antes de calcular, con aritm\'etica visible y acciones de mejora implementadas y re-medidas.
  \item Verificar y cerrar la trazabilidad end-to-end en la cadena Fuente $\rightarrow$ RF $\rightarrow$ CU $\rightarrow$ Clase $\rightarrow$ Estado $\rightarrow$ BDD $\rightarrow$ CP, eliminando requisitos hu\'erfanos y cadenas rotas.
  \item Especificar al menos dos componentes de inteligencia artificial con sus bloques de descripci\'on, datos, m\'etricas de \'exito, \'etica/privacidad y monitoreo, incluyendo requisitos de equidad y explicabilidad.
  \item Consolidar la versi\'on final del ERS (v2.0) y el informe integrador de las cinco unidades con referencias IEEE verificables.
  \item Defender individualmente ante el tribunal cualquier decisi\'on de ingenier\'ia tomada en el proyecto, anclando cada respuesta en el artefacto correspondiente.
  \item Ejecutar una retrospectiva de equipo y documentar el aporte individual verificable de cada integrante contra el historial del repositorio.
\end{enumerate}

\subsection{Alcance del documento y regla de autenticidad}

Este informe cubre el ciclo completo de ingenier\'ia de requisitos ejecutado en cinco pr\'acticas experimentales (PE1--PE5): elicitaci\'on multi-t\'ecnica con evidencia marcada (EV-01 a EV-07, entrevistas ENTR-01 a ENTR-06), an\'alisis y codificaci\'on tem\'atica, modelado organizacional y UML, especificaci\'on conforme a ISO/IEC/IEEE 29148:2018, validaci\'n por inspecci\'on con registro de defectos y re-inspecci\'on, auditor\'ia m\'etrica final, trazabilidad end-to-end, requisitos de IA y preparaci\'on de la defensa t\'ecnica \cite{iso29148}. No cubre el dise\~no detallado, la implementaci\'on ni el entrenamiento efectivo de los modelos de IA: lo que aqu\'i se entrega es la especificaci\'on verificable que sustentar\'a esas fases posteriores.

Siguiendo la regla de autenticidad de la pr\'actica, toda cifra reportada en este documento proviene de artefactos contables del repositorio del equipo: conteos base versionados antes de calcular, scripts PowerShell reproducibles con verificación por hash SHA-256, matriz de trazabilidad fila a fila y registro de defectos fechado. Ninguna m\'etrica se reporta sin el conteo que la respalda \cite{wohlin2012}.

\subsection{Estructura del documento}

El resto del informe se organiza as\'i: la Secci\'on~2 describe la metodolog\'ia de ingenier\'ia de requisitos seguida en PE1--PE5 y justifica las decisiones metodol\'ogicas; la Secci\'on~3 presenta la ERS/SRS final v2.0 y su evoluci\'on de l\'inea base; la Secci\'on~4 interpreta los modelos UML principales; la Secci\'on~5 documenta la validaci\'n (inspecci\'on, correcciones y re-inspecci\'on); la Secci\'on~6 integra la gesti\'n del proyecto y la matriz de trazabilidad end-to-end; la Secci\'on~7 especifica los requisitos de los componentes de inteligencia artificial; la Secci\'on~8 reporta la auditor\'ia de calidad con las seis m\'etricas ISO/IEC 25010:2023 y su comparaci\'n antes/despu\'es; la Secci\'on~9 recoge la retrospectiva Start--Stop--Continue; y la Secci\'on~10 cierra con cinco conclusiones integradoras, una por unidad. Los anexos A--E contienen el instrumento de auditor\'ia, el banco de preguntas del tribunal con respuestas preparadas, la declaraci\'n de aporte individual, la retrospectiva tabulada y la declaraci\'n de uso de asistentes de IA, respectivamente.
